\clearpage
\phantomsection
\addcontentsline{toc}{chapter}{Liste des acronymes}
\thispagestyle{plain}
\setlength{\parskip}{0pt}
\setlength{\baselineskip}{12pt}

\begin{center}
{\Large\bfseries LISTE DES ACRONYMES}
\end{center}
\vspace{0.4em}

{\footnotesize
\renewcommand{\arraystretch}{1.05}
\noindent
\begin{tabular}{@{}>{\bfseries}l@{\ :\ }p{0.78\textwidth}@{}}
ACO & Ant Colony Optimization (optimisation par colonies de fourmis) \\
BCD & Boustrophedon Cellular Decomposition (décomposition cellulaire boustrophédon) \\
CDPM & Cable-Driven Parallel Manipulator (manipulateur parallèle à câbles) \\
CDPR & Cable-Driven Parallel Robot (robot parallèle à câbles) \\
CPP & Coverage Path Planning (planification de chemin de couverture complète) \\
CR & Cable Robot (robot à câbles) \\
CSI & Camera Serial Interface (interface série pour caméra) \\
CSV & Comma-Separated Values (valeurs séparées par des virgules) \\
DDL & Degré(s) de liberté \\
FOV & Field Of View (champ de vision) \\
GA & Genetic Algorithm (algorithme génétique) \\
LP & Linear Programming (programmation linéaire) \\
LSPB & Linear Segments with Parabolic Blends (segments linéaires avec raccords paraboliques) \\
MGD & Modèle Géométrique Direct \\
MGI & Modèle Géométrique Inverse \\
MOEA/D & Multi-Objective Evolutionary Algorithm based on Decomposition \\
 & (algorithme évolutionnaire multi-objectif par décomposition) \\
MOPSO & Multi-Objective Particle Swarm Optimization (optimisation multi-objectif par essaim particulaire) \\
NEMA & National Electrical Manufacturers Association \\
NP & Nondeterministic Polynomial time (polynomial non déterministe) \\
NSGA-II & Non-dominated Sorting Genetic Algorithm II (algorithme génétique à tri non dominé II) \\
OX & Order Crossover (croisement d'ordre) \\
PM & Plateforme Mobile \\
PMX & Partially Mapped Crossover (croisement partiellement apparié) \\
PRM & Probabilistic Roadmap (feuille de route probabiliste) \\
PSO & Particle Swarm Optimization (optimisation par essaim particulaire) \\
QP & Quadratic Programming (programmation quadratique) \\
RGB & Red, Green, Blue (rouge, vert, bleu) \\
RRT & Rapidly-exploring Random Tree (arbre aléatoire à exploration rapide) \\
SBX & Simulated Binary Crossover (croisement binaire simulé) \\
SEW & Static Equilibrium Workspace (espace d'équilibre statique) \\
SLSQP & Sequential Least Squares Programming (programmation séquentielle par moindres carrés) \\
TSP & Traveling Salesman Problem (problème du voyageur de commerce) \\
VRP & Vehicle Routing Problem (problème de tournées de véhicules) \\
VRPTW & Vehicle Routing Problem with Time Windows (tournées de véhicules avec fenêtres temporelles) \\
WCW & Wrench Closure Workspace (espace de fermeture en effort) \\
WFW & Wrench-Feasible Workspace (espace faisable en effort) \\
\end{tabular}
}

